\documentclass[dvipdfmx]{beamer}
%テーマ
\usetheme[numbering=counter,shape=square]{LightTheme}
%パッケージ
\usepackage{pxjahyper}
\usepackage{tikz}
%tikzライブラリ
\usetikzlibrary{arrows.meta}
\usetikzlibrary{positioning}
\usetikzlibrary{shapes.callouts}
%フォント
\renewcommand{\kanjifamilydefault}{\gtdefault}

%タイトル設定
\title{TechRingにかける想い}
\subtitle{学科を超えて}
\author[Jin]{石田　優仁}
\institute[TechRing]{奈良高専TechRing}
\date[2023/04/30]{2023/04/30}

\begin{document}

\maketitle

\begin{frame}{高専どんなところ？}

  \begin{block}{国立高等専門学校の特徴}
    \begin{itemize}
      \setlength{\itemsep}{1mm}
      \item 5年制、一貫のカリキュラム
      \item 電気電子系や機械系に限らず、物質や建築、商船もある
      \item 高専祭という文化祭での文化発表が活発
      \item ロボコンやデザコン、プロコンなどのコンテストが充実
      \item 高専カンファレンスや様々な技術コミュニティでの交流
    \end{itemize}
  \end{block}

  \begin{alertblock}{奈良高専の特徴}
      \begin{itemize}
        \setlength{\itemsep}{1mm}
        \item 来年で創設60周年を迎える
        \item 機械、電気、電子制御、情報、物質科学の5学科
        \item 最寄り駅からバスで20分ほど
        \item 女子枠推薦がある
      \end{itemize}
  \end{alertblock}

\end{frame}

\begin{frame}{高専スタートアップ事業について}{時代は大高専時代に...}

  \begin{alertblock}{高専スタートアップ事業概要}
      \begin{itemize}
        \setlength{\itemsep}{1mm}
        \item アントレプレナーシップ教育に取り組究極の
        国公私立高専の戦略的な取組に必要な経費を補助
        \item 高専生の活動を後押しし\\
        学生が自らの技術を用いた社会課題解決に取り組み
        \item 予算は57高専で合計\alert{60億円}も
      \end{itemize}
  \end{alertblock}

  \begin{center}
    \begin{tikzpicture}
      \draw[-{Latex[length=5mm]},line width = 3pt,red] (0,.7) -- (0,0);
    \end{tikzpicture}
  \end{center}

  \vspace{-3.5mm}

  \begin{greyblock}{スタートアップ事業STEP}
    \begin{enumerate}
      \setlength{\itemsep}{1mm}
      \item アントレプレナーシップや企業について知る
      \item 高専生のチャレンジを応援するための環境整備\\
      $\Rightarrow$\alert{起業家工房}の整備
      \item 地域社会の課題解決に挑戦
    \end{enumerate}
  \end{greyblock}
\end{frame}

\begin{frame}{TechRingとは？}
    \begin{description}[一つの輪に]
      \item[学科の壁] 各学科には様々な長所や短所が混在するが...\par
      各学科が壁で隔てられ、\\長所を活かすことも短所を補うこともできない！\par
      \hspace{0.3\textwidth}
      \begin{tikzpicture}
        \draw[-{Latex[length=5mm]},line width = 3pt,red] (0,.7) -- (0,0);
      \end{tikzpicture}
      \vspace{-1.5mm}
      \item[一つの輪に] 学科を超えた相互補完型の交流で\\究極のアントレプレナーシップを目指す！
    \end{description}

\end{frame}

\begin{frame}{TechRingの目的}
  TechRingの目指す目的を紹介
\end{frame}

\begin{frame}{TechRingの担う学内活動 その１}{イベント開催}
  奈良高専生に向けてイベントを開催する話\\
  LTイベントや何か製作イベントなど
\end{frame}

\begin{frame}{TechRingの担う学内活動 その２}{技術交流}
  5団体でノウハウや活動の共有をしたり、\\
  コンテストの共同参加やその他技術交流に関する話
\end{frame}

\begin{frame}{TechRingの担う学内活動 その３}{共同制作}
  技術交流を通して共同で制作をする話\\
  共同制作したものを外部公開する話
\end{frame}

\begin{frame}{TechRingを創る仲間たち その１}{機械研究会}
  機械研究会にフォーカス\\
  活動の紹介と出展の紹介を
\end{frame}

\begin{frame}{TechRingを創る仲間たち その２}{MeCafe}
  MeCafeにフォーカス\\
  活動の紹介と出展の紹介を
\end{frame}

\begin{frame}{TechRingを創る仲間たち その３}{電気技術研究会}
  電気技術研究会にフォーカス\\
  活動の紹介と出展の紹介を
\end{frame}

\begin{frame}{TechRingを創る仲間たち その４}{システム開発研究会}
  システム開発研究会にフォーカス\\
  活動の紹介
\end{frame}

\begin{frame}{TechRingを創る仲間たち その５}{情報処理研究会}
  情報処理研究会にフォーカス\\
  活動の紹介と出展の紹介を
\end{frame}

\begin{frame}{TechRingのこれから}
  TechRingの今後の活動を紹介
\end{frame}

\begin{frame}
  \begin{Large}
    ご清聴、ありがとうございました。
  \end{Large}
\end{frame}

\end{document}